% Claim proof file for row claim_3_17 -- "MarginalizationsCommute".
%
% Proof lifted from lecture-notes/lecture_notes/graphs.tex:1007--1117 (the
% \Claude{...} block immediately following the claim). Light edits only:
% an explicit \refrow{def_3_14} cross-link the first time the marginalization
% 4-tuple is unpacked, and the LN's outer \Claude{...} wrapper stripped.
% The %\TODO{...} comment-note at LN line 1110 is preserved verbatim (as a
% TeX `%` comment) for audit-trail purposes; it does not affect compilation.
%
% Compiles standalone via the `subfiles` package.

\documentclass[main]{subfiles}

\begin{document}
% ref: claim_3_17 (proof, with statement restated for readability)
% title: MarginalizationsCommute
\def\rowref{claim\_3\_17}
\phantomsection\label{claim_3_17}
%
% Cross-references: see claim_<ref>_statement_<title>.tex in this folder
% for the corresponding statement.

% --- statement (restated) ---------------------------------------------------
\begin{Lem}[Marginalizations commute]\label{marginalizations-commute}
      Let $G=(J,V,E,L)$ be a CDMG and $W_1, W_2 \ins V$ two disjoint subsets of output nodes.
      Then we have:
      \[ \lp G^{\sm W_1} \rp^{\sm W_2} = \lp G^{\sm W_2} \rp^{\sm W_1} =  G^{\sm (W_1 \cup W_2)}. \]
     % This shows that the notation:
     % \[ G(V \sm (W_1 \cup W_2)|\doit(J))  \]
     % is unambiguous.
\end{Lem}

% --- proof ------------------------------------------------------------------
\begin{proof}
We show $\lp G^{\sm W_1}\rp^{\sm W_2} = G^{\sm(W_1 \cup W_2)}$; the other equality then follows by symmetry (swapping $W_1$ and $W_2$).
By \refrow{def_3_14} each side is a CDMG presented as a 4-tuple $(J_\ast, V_\ast, E_\ast, L_\ast)$.
Both graphs have input nodes $J$ and output nodes $V \sm (W_1 \cup W_2)$.
It remains to show that the directed and bidirected edge sets coincide.

\medskip\noindent\emph{Directed edges.}
We show: for $\ul{v},\ol{v} \in J \cup V \sm (W_1 \cup W_2)$,
\[
  \ul{v} \tuh \ol{v} \in (E^{\sm W_1})^{\sm W_2}
  \quad\iff\quad
  \ul{v} \tuh \ol{v} \in E^{\sm(W_1 \cup W_2)}.
\]

$(\Rightarrow)$
If $\ul{v} \tuh \ol{v} \in (E^{\sm W_1})^{\sm W_2}$, there is a directed walk
$\ul{v} = u_0 \tuh u_1 \tuh \cdots \tuh u_m = \ol{v}$
in $G^{\sm W_1}$ with all intermediate nodes $u_1,\dots,u_{m-1} \in W_2$.
Each edge $u_i \tuh u_{i+1} \in E^{\sm W_1}$ corresponds to a directed walk
in $G$ from $u_i$ to $u_{i+1}$ with intermediate nodes in $W_1$.
Concatenating all these walks yields a directed walk from $\ul{v}$ to $\ol{v}$
in $G$ with all intermediate nodes in $W_1 \cup W_2$, so
$\ul{v} \tuh \ol{v} \in E^{\sm(W_1 \cup W_2)}$.

$(\Leftarrow)$
If $\ul{v} \tuh \ol{v} \in E^{\sm(W_1 \cup W_2)}$, there is a directed walk
$\ul{v} = a_0 \tuh a_1 \tuh \cdots \tuh a_n = \ol{v}$
in $G$ with all intermediate nodes $a_1,\dots,a_{n-1} \in W_1 \cup W_2$.
Extract the subsequence of nodes not in $W_1$:
$\ul{v}=b_0, b_1, \dots, b_r, b_{r+1}=\ol{v}$ with $b_1,\dots,b_r \in W_2$.
Between consecutive nodes $b_j$ and $b_{j+1}$, all intermediate nodes of
the original walk lie in $W_1$, so $b_j \tuh b_{j+1} \in E^{\sm W_1}$.
This gives a directed walk from $\ul{v}$ to $\ol{v}$ in $G^{\sm W_1}$
with intermediate nodes in $W_2$, hence
$\ul{v} \tuh \ol{v} \in (E^{\sm W_1})^{\sm W_2}$.

\medskip\noindent\emph{Bidirected edges.}
We show: for distinct $\ul{v}, \ol{v} \in V \sm (W_1 \cup W_2)$,
\[
  \ul{v} \huh \ol{v} \in (L^{\sm W_1})^{\sm W_2}
  \quad\iff\quad
  \ul{v} \huh \ol{v} \in L^{\sm(W_1 \cup W_2)}.
\]

$(\Rightarrow)$
If $\ul{v} \huh \ol{v} \in (L^{\sm W_1})^{\sm W_2}$, there exists a bifurcation
in $G^{\sm W_1}$ between $\ul{v}$ and $\ol{v}$ with all intermediate nodes in $W_2$:
\[
  \ul{v} \hut u_1 \hut \cdots \hut u_{k-1} \hus u_k \tuh \cdots \tuh u_{m-1} \tuh \ol{v}
\]
with $u_1,\dots,u_{m-1} \in W_2$.
Each directed edge on the bifurcation (in $E^{\sm W_1}$) expands to a directed walk
in $G$ with intermediate nodes in $W_1$.
If the hinge edge $u_{k-1} \hus u_k$ is directed ($u_{k-1} \hut u_k \in E^{\sm W_1}$),
it likewise expands to a directed walk in $G$ through $W_1$.
If the hinge edge is bidirected ($u_{k-1} \huh u_k \in L^{\sm W_1}$),
it expands to a bifurcation in $G$ between $u_{k-1}$ and $u_k$ with intermediate nodes in $W_1$.
Substituting all these expansions into the bifurcation yields a bifurcation in $G$
between $\ul{v}$ and $\ol{v}$ with all intermediate nodes in $W_1 \cup W_2$,
so $\ul{v} \huh \ol{v} \in L^{\sm(W_1 \cup W_2)}$.

$(\Leftarrow)$
If $\ul{v} \huh \ol{v} \in L^{\sm(W_1 \cup W_2)}$, there exists a bifurcation in $G$:
\[
  \ul{v} = a_0 \hut a_1 \hut \cdots \hut a_{k-1} \hus a_k \tuh \cdots \tuh a_{n-1} \tuh a_n = \ol{v}
\]
with all intermediate nodes $a_1,\dots,a_{n-1} \in W_1 \cup W_2$.
We construct a bifurcation in $G^{\sm W_1}$ between $\ul{v}$ and $\ol{v}$ with intermediate nodes in $W_2$.

\emph{Left arm.}
Reading the left arm in forward direction gives a directed walk in $G$:
$a_{k^*} \tuh \cdots \tuh a_1 \tuh \ul{v},$
where $a_{k^*}$ equals $a_k$ if the hinge is directed ($a_{k-1} \hut a_k$) and
$a_{k^*}$ equals $a_{k-1}$ if the hinge is bidirected ($a_{k-1} \huh a_k$).
All intermediate nodes of this directed walk lie in $W_1 \cup W_2$.
Extracting the subsequence of non-$W_1$ nodes and applying the same argument as
for directed edges gives a directed walk from $a_{k^*}$ (or its
first non-$W_1$ descendant on this arm) to $\ul{v}$ in $G^{\sm W_1}$ with
intermediate nodes in $W_2$.

\emph{Right arm.}
Similarly, $a_k \tuh \cdots \tuh a_{n-1} \tuh \ol{v}$ is a directed walk in $G$
with intermediate nodes in $W_1 \cup W_2$.
The non-$W_1$ extraction gives a directed walk in $G^{\sm W_1}$ from
$a_k$ (or its first non-$W_1$ descendant on this arm) to $\ol{v}$
with intermediate nodes in $W_2$.

\emph{Hinge.}
Let $\beta$ be the first non-$W_1$ node reached from the hinge going along
the left arm towards $\ul{v}$ (or $\ul{v}$ itself if all intermediate left-arm
nodes lie in $W_1$).
Let $\gamma$ be the first non-$W_1$ node reached from the hinge going along
the right arm towards $\ol{v}$ (or $\ol{v}$ itself if all intermediate right-arm
nodes lie in $W_1$).

If $\beta \neq \gamma$: the subwalk from $\beta$ backwards along the left arm,
through the hinge, and forward along the right arm to $\gamma$
is a bifurcation in $G$ with all intermediate nodes in $W_1$, so
$\beta \huh \gamma \in L^{\sm W_1}$.

If $\beta = \gamma$: all nodes between $\beta$ and the hinge on both arms lie in $W_1$,
so the node $\beta$ is the source of directed walks to both $\ul{v}$ and $\ol{v}$
in $G^{\sm W_1}$ (with intermediate nodes in $W_2$).
%\TODO{The previous version incorrectly claimed ``the case $\beta=\gamma$ with a bidirected hinge cannot occur'', but the at-most-one-arrowhead condition is per occurrence, not per node, so this case CAN occur. The conclusion is still correct: $\beta=\gamma$ is a source of directed walks to both endpoints regardless of hinge type.}

In all cases, assembling the contracted left arm, the contracted hinge
($\beta \huh \gamma$ or source~$\beta{=}\gamma$), and the contracted right arm yields
a bifurcation in $G^{\sm W_1}$ between $\ul{v}$ and $\ol{v}$ with all intermediate
nodes in $W_2$.
Hence $\ul{v} \huh \ol{v} \in (L^{\sm W_1})^{\sm W_2}$.
\end{proof}

\end{document}
